\section{Equations of motion}\label{app:eoms}

The action in the gravity dual consists of Einstein plus ordinary and ghost scalars
\begin{equation}
	S_{d+1}=\frac{1}{16\pi G_{d+1}}\int \dd^{d+1} x\, \sqrt{-G}\left[R+\frac{1}{2}\partial_M \eta^a\partial^M\eta^a-\frac{1}{2}\partial_M \phi^i\partial^M\phi^i-V\left(\{\eta^a,\phi^i\}\right)\right]\ .
\end{equation}
Einstein's equations are
\begin{equation}
	R_{MN}=\frac{1}{2}\widetilde T_{MN}\equiv \frac{1}{2}\left(T_{MN}-\frac{1}{d-1}G_{MN}T^L_{\ L} \right)\ ,
\end{equation}
where
\begin{equation}
	T_{MN}=-\partial_M \eta^a\partial_N\eta^a+\partial_M \phi^i\partial_N\phi^i-\frac{1}{2}G_{MN}\left( -\partial_L \eta^a\partial^L\eta^a+\partial_L \phi^i\partial^L\phi^i+2V\right)\ ,
\end{equation}
hence
\begin{equation}
	\widetilde T_{MN}=-\partial_M \eta^a\partial_N\eta^a+\partial_M \phi^i\partial_N\phi^i+\frac{2}{d-1}G_{MN}V\ .
\end{equation}
The scalar equations are
\begin{equation}
	\nabla^M\partial_M\phi^i-\partial_{\phi^i}V=0
	\ ,\qquad \nabla^M\partial_M\eta^a+\partial_{\phi^a}V=0\ .
\end{equation}

\subsection{Equations for the background}

The background metric and scalars for a domain wall solution are 
\begin{equation}
	\dd s^2=g_{MN}\dd x^M \dd x^N=\dd r^2+e^{2A(r)}\eta_{\mu\nu}\dd x^\mu \dd x^\nu
	\ ,\qquad 
	\phi^i=\phi^i_0(r)
	\ ,\qquad 
	\eta^a=\eta^a(r)\ .
\end{equation}
Einstein's equations can be simplified to
\begin{equation}\label{eq:Einsbackg}
	A''+d(A')^2+\frac{V_0}{d-1}=0
	\ ,\qquad 
	(d-1)A''=\frac{1}{2}({\eta_0^a}')^2-\frac{1}{2}({\phi_0^i}')^2\ .
\end{equation}
Where $V_0=V(\{\eta^a_0,\phi^i_0\})$. The equations for the background scalars are
\begin{equation}\label{eq:scalarbackg}
	{\phi_0^i}''+dA'{\phi_0^i}'-\partial_{\phi^i}V_0=0
	\ ,\qquad {\eta_0^a}''+dA'{\eta_0^a}'+\partial_{\eta^a}V_0=0\ .
\end{equation}
One can check that solutions to the first order equations satisfy the second order equations as well
\begin{equation}
   A'=-\frac{W_0}{2(d-1)}
   \ ,\qquad {\phi^i}'=\partial_{\phi^i}W_0
   \ ,\qquad
   {\eta^a}'=-\partial_{\eta^a}W_0\ ,
\end{equation}
where $W_0=W(\{\eta^a_0,\phi^i_0\})$ and
\begin{equation}
	V=\frac{1}{2}(\partial_{\phi^i}W)^2-\frac{1}{2}(\partial_{\eta^a} W)^2-\frac{d}{4(d-1)}W^2\ .
\end{equation}

\subsection{Linearized fluctuations}

We will expand to linear order in the metric and the scalar fields around the domain wall background
\begin{equation}
	G_{MN}=g_{MN}+h_{MN}
	\ ,\qquad \phi^i=\phi^i_0+\varphi^i
	\ ,\qquad 
	\eta^a=\eta^a_0+\chi^a\ .
\end{equation}
The linear order Ricci tensor has the form
\begin{equation}
   \delta R_{MN}= -\frac{1}{2}\left(\nabla^2 h_{MN}-2\nabla^L \nabla_{(M}h_{N)L}+\nabla_M \partial_N h\right)\ .
\end{equation}
where covariant derivatives are with respect to the background metric and indices are raised and lowered with the background metric $h=g^{MN}h_{MN}$.\\
The energy-momentum tensor of the scalars to linear order is
\begin{equation}
\begin{split}
     \delta \widetilde T_{MN}=&-2\partial_{(M}\eta_0^a \partial_{N)}\chi^a+2\partial_{(M}\phi_0^i \partial_{N)}\varphi^i\\
     &+\frac{2}{d-1}h_{MN}V_0+\frac{2}{d-1}g_{MN}\partial_{\phi^i} V_0\,\varphi^i+\frac{2}{d-1}g_{MN}\partial_{\eta^a} V_0\,\eta^a\ .
	\end{split}
\end{equation}
Expanding to linear order the scalar equations of motion, we get
\begin{eqnarray}
\notag \nabla^M \partial_M \varphi^i-\partial_{\phi^i}\partial_{\phi^j} V_0\,\varphi^j-\partial_{\phi^i}\partial_{\eta^a} V_0\,\chi^a-h^{MN}\nabla_M\partial_N \phi_0^i-\left(\nabla_M h^{MN}-\frac{1}{2}\partial^N h\right)\partial_N\phi^i_0 &=&0, \\
& & \\
\notag \nabla^M \partial_M \chi^a+\partial_{\eta^a}\partial_{\eta^b} V_0\,\chi^b+\partial_{\eta^a}\partial_{\phi^i} V_0\,\varphi^i-h^{MN}\nabla_M\partial_N \eta_0^a-\left(\nabla_M h^{MN}-\frac{1}{2}\partial^N h\right)\partial_N\eta^a_0 &=&0.\\ & &
\end{eqnarray}
In the domain wall background the only non-zero Christoffel symbols are
\begin{equation}
	\Gamma^r_{\mu\nu}=-A' e^{2A}\eta_{\mu\nu}=-A' g_{\mu\nu}
	\ , \qquad 
	\Gamma^\mu_{r\nu}=A' \delta^\mu_\nu\ .
\end{equation}
Then, the scalar equations become, after using the background equations of motion
\begin{eqnarray}
\notag \partial^M \partial_M \varphi^i+dA'\partial_r\varphi^i -\partial_{\phi^i}\partial_{\phi^j} V_0\,\varphi^j-\partial_{\phi^i}\partial_{\eta^a} V_0\,\chi^a-\partial_{\phi^i} V_0 h_{rr} -\left(\partial^M h_{Mr}-\frac{1}{2}\partial_r h\right){\phi_0^i}' &=&0, \\ & & \\
\notag \partial^M \partial_M \chi^a+dA'\partial_r\chi^a +\partial_{\eta^a}\partial_{\eta^b} V_0\,\chi^b+\partial_{\eta^a}\partial_{\phi^i} V_0\,\varphi^i+\partial_{\eta^a} V_0 h_{rr} -\left(\partial^M h_{Mr}-\frac{1}{2}\partial_r h\right){\eta_0^a}' &=&0 .\\ \label{eq:chieq} & &
\end{eqnarray}
The components of the energy-momentum tensor are
\begin{eqnarray}
	\delta \widetilde T_{rr} &=& -2{\eta_0^a}'\partial_r \chi^a+2{\phi_0^i}'\partial_r \varphi^i+\frac{2}{d-1}\left(\partial_{\phi^i} V_0 \varphi^i+\partial_{\eta^a} V_0 \chi^a+ V_0 h_{rr}\right),\\
	\delta \widetilde T_{\mu r} &=& -{\eta_0^a}' \partial_\mu\chi^a+{\phi^i_0}'\partial_\mu \varphi^i+\frac{2}{d-1}V_0 h_{\mu r},\\
	\delta \widetilde T_{\mu\nu} &=& \frac{2}{d-1}g_{\mu\nu} \left(\partial_{\phi^i} V_0\, \varphi^i+\partial_{\eta^a} V_0\, \chi^a\right)+\frac{2}{d-1}V_0 h_{\mu\nu}.
\end{eqnarray}
The components of the Ricci tensor can be written as
\begin{eqnarray}
 \notag   -2\delta R_{rr} &=&\partial^M\partial_M h_{rr}-2\partial_r\left(\partial^M h_{Mr}-\frac{1}{2}\partial_rh\right)-4 A'\left(\partial^M h_{Mr} -\frac{1}{2}\partial_r h\right)\\
	& & -(d-2)A' \partial_r h_{rr}\ .\\
 \notag    -2\delta R_{r\mu} &=& \partial^M\partial_M h_{r\mu}-\partial_r\left(\partial^M h_{M\mu}-\frac{1}{2}\partial_\mu h\right)-\partial_\mu\left(\partial^M h_{Mr}-\frac{1}{2}\partial_rh\right)+2 A'' h_{\mu r}
 \\
 & & + 2d(A')^2 h_{\mu r}-(d-1) A'\partial_\mu h_{rr}
 \ ,\\
\notag    -2\delta R_{\mu\nu} &=& \partial^M\partial_M h_{\mu\nu}-2\partial_{(\mu}\left(\partial^M  h_{\nu) M}-\frac{1}{2}\partial_{\nu)} h\right)-2 A'' g_{\mu\nu}h_{r r}\\
\notag & & +(A')^2\left(4 h_{\mu\nu}-2d g_{\mu\nu} h_{rr}\right)\\
 & & +A'\left[-2(d-2)\partial_{(\mu} h_{\nu)r}-2g_{\mu\nu} \left(\partial^M h_{Mr}-\frac{1}{2}\partial_r h\right)+(d-4)\partial_r h_{\mu\nu}\right]\ .
\end{eqnarray}
In the case of a metric vector fluctuation with $\varphi^i=0$, $\eta^a=0$, $h_{rr}=0$ and $h_{\mu\nu}=0$, the scalar equations and the Einstein equation for the $rr$ components become
\begin{equation}\label{eq:transvapp}
	\partial^\mu h_{\mu r}=0\ .
\end{equation}
The Einstein equations for the $r\mu$ components become
\begin{equation}
\begin{split}
0=& -2\delta R_{r\mu}+\delta \widetilde T_{r\mu}\\ &=\partial^M\partial_M h_{\mu r}-\partial_r^2 h_{r\mu}+2\left[A''+d(A')^2+\frac{1}{d-1} V_0\right] h_{r\mu}=g^{\alpha\beta}\partial_\alpha \partial_\beta h_{r\mu}\ .
\end{split}
\end{equation}
The equation for the $\mu\nu$ components is
\begin{align}
0=& -2\delta R_{\mu\nu}+\delta \widetilde T_{\mu\nu}
=-2\partial_r\left(\partial_{(\mu}h_{\nu)r}\right)-2(d-2)A'\partial_{(\mu}h_{\nu)r}-2A' g_{\mu\nu}\partial^\alpha h_{\alpha r}
\\ \nonumber&=-2e^{-(d-2)A}\partial_{(\mu}\partial_r\left( e^{(d-2)A} h_{\nu)r}\right).
\end{align}
Where we used \eqref{eq:transvapp} in the last step.

A scalar fluctuation corresponds to the ansatz
\begin{equation}
	h_{r\mu}=\partial_\mu h_V
	\ ,\qquad h_{\mu\nu}=e^{2A}\eta_{\mu\nu}h_S+\partial_\mu\partial_\nu h_L\ .
\end{equation}
A massless fluctuation would obey the equations
\begin{subequations}\label{eq:dilatonmode}
\begin{eqnarray}
	(r\mu): \ 0 &=& \partial_r h_S-A' h_R-\frac{\eta_0'}{d-1}\chi
	\ ,\\
	(rr): \ 0&=& \partial_r^2 h_S-A'\partial_r h_R-\frac{2}{d-1}\eta_0'\partial_r \chi
	\ ,\\
	\eta^{\mu\nu}(\mu\nu): \ 0&=&(d-2) A'' h_R+(\eta_0''\chi-\eta_0' \partial_r\chi)\\
	0&=& \partial_r^2 \chi+dA'\partial_r\chi+\partial_\eta^2 V\chi+\partial_\eta V h_R-\frac{1}{2}(\partial_r h_R-d\partial_r h_S)\eta_0'\ .
\end{eqnarray}
\end{subequations}
And there is an additional equation $\sim \partial_\mu\partial_\nu$ involving $h_V$ and $h_L$.
The difference of the first two equations gives
\begin{equation}
	(d-1)A''h_R+(\eta_0''\chi-\eta_0'\partial_r\chi)=0\ .
\end{equation}
Thus, we get
\begin{equation}\label{eq:dilatonsolution}
	h_R=0
	\ ,\qquad \chi=\tau(x) \eta_0'(r)
	\ ,\qquad
	\partial_r h_S=\frac{(\eta_0')^2}{d-1}\tau(x)\ .
\end{equation}
One can check that the last equation in \eqref{eq:dilatonmode} is satisfied using the first order equations and the definition of the potential in terms of the superpotential. Usually this would be the dilaton mode, but, if we integrate the equation for $h_S$ in \eqref{eq:dilatonsolution}, we cannot fix the integration constant in such a way that $h_S$ vanishes at both boundaries. Thus, the would-be dilaton mode cannot be made normalizable at both boundaries simultaneously.

\subsubsection{Equations in radial gauge for a single ghost field}

In order to study the stability of the solutions when there is a single ghost field we will work in the radial gauge $h_{M r}=0$ and expand in plane waves along the field theory directions
\begin{equation}
\begin{split}
	&h_{\mu\nu}(r,x)=\int \frac{\dd^dk}{(2\pi)^d} e^{ik\cdot x}H_{\mu\nu}(r,k),\\
	&\chi(r,x)=\int \frac{\dd^dk}{(2\pi)^d} e^{ik\cdot x}X(r,k).
\end{split}
\end{equation}
In the following we will use $\partial_r=\,'$ for the Fourier components, as we did for the background solutions. 

The equations of motion are invariant under linearized diffeomorphisms
\begin{equation}
	\delta h_{MN}=-g_{ML}\partial_N \xi^L-g_{NL}\partial_M \xi^L-\xi^L\partial_L g_{MN},\quad \delta \chi=-\xi^L\partial_L\eta_0.
\end{equation}
Fixing the radial gauge does not completely fix all diffeomorphisms, there is a remnant which takes the form
\begin{equation}
	\xi^r(r,x)=\xi^r(x),\quad \xi^\mu(r,x)=\eta^{\mu\nu}\left(\lambda_\mu(x)-\partial_\nu\xi^r(x)\int_{r_0}^r \dd u\, e^{-2A(u)}\right).
\end{equation}
The transformation of the fluctuations under the remnant diffeomorphism is
\begin{equation}
	\delta \chi=-\xi^r\eta_0'
	\ ,\quad 
	\delta h_{\mu\nu}=-2e^{2A}\partial_{(\mu}\lambda_{\nu)}-2 A' e^{2A}\eta_{\mu\nu}\xi^r-\partial_\mu\partial_\nu\xi^r\int_{r_0}^r \dd u\, e^{-2A(u)}\ .
\end{equation}
We can construct invariant combinations adding some of the fluctuations with appropriate factors. In order to simplify the analysis, let us work with the modes in the plane wave expansion, and introduce the following $d-1$ transverse polarization vectors with respect to the (timelike) momentum $k^\mu=(\omega,\bm{k})$,
\begin{equation}
q_\mu=\frac{1}{\sqrt{\omega^2-\bm{k}^2}}\left(|\bm{k}|\ ,-\omega\frac{\bm{k}}{|\bm{k}|}\right)\ ,\quad e_\mu^\alpha=(0,\bm{e}^\alpha), \ \alpha=1,\dots,d-2\ ,
\end{equation}
where
\begin{equation}
     \bm{k}\cdot\bm{e}^\alpha=0
     \ ,\qquad \bm{e}^\alpha\cdot\bm{e}^\beta=\delta^{\alpha\beta}\ .
\end{equation}
In the following, we will use the Minkowski metric to raise and lower indices along the field theory directions. There are two independent scalar invariant combinations
\begin{eqnarray}
 \label{eq:zetaeta}   
 Z_\eta=X-\frac{\eta_0'e^{-2A}}{2A'}\frac{1}{d-2}\sum_{\alpha=1}^{d-2}e_\mu^\alpha e_\nu^\alpha H^{\mu\nu}
 ,\\
 \label{eq:HS}
 H_S=\left(q_\mu q_\nu-\frac{1}{d-2}\sum_{\alpha=1}^{d-2}e_\mu^\alpha e_\nu^\alpha\right)H^{\mu\nu}.
\end{eqnarray}
There are $d-2$ invariant vector modes 
\begin{equation}
	H_V^\alpha=q_\mu e_\nu^\alpha H^{\mu\nu}.
\end{equation}
Finally, the invariant tensor modes are, for $\alpha< \beta$,
\begin{equation}
\begin{split}
&H_1^{\alpha\beta}=H_1^{\beta\alpha}=e^{\alpha}_\mu e^{\beta}_\nu H^{\mu\nu},\\
&H_2^{\alpha\beta}=-H_2^{\beta\alpha}=(e_\mu^{\alpha}e_\nu^{\alpha}-e_\mu^{\beta}e_\nu^{\beta})H^{\mu\nu}.
\end{split}
\end{equation}
This gives $\frac{(d-2)(d-1)}{2}-1$ independent tensor modes. In total there are $1+\frac{(d-2)(d+1)}{2}$
invariant modes, that agree with the $1+\frac{d(d+1)}{2}$ fluctuations of the metric and scalar, minus the $d+1$ constraints that we obtain fixing the radial gauge. The constraints remove the $d$ longitudinal components of the metric and the trace part of the transverse components.

We can rearrange the equations of the linearized fluctuations to obtain a set of minimally coupled equations for the invariant modes in the radial gauge. The equations for the invariant modes not involving the scalar field are found simply by contracting the $\mu\nu$ Einstein equations with the projectors that appear in their definitions. Taking also into account that these modes are transverse and traceless, one gets that each of them is decoupled from the rest of the modes and satisfies the same equations ($H=H_S, H_V^\alpha,H_1^{\alpha\beta},H_2^{\alpha\beta}$)
\begin{equation}
	H''+(d-4)A' H'+\left[e^{-2A}M^2+4(A')^2+\frac{2}{d-1}V_0\right]H=0\ ,
\end{equation}
where we have defined $k^2=k^\mu k_\mu=-M^2$. We can combine \eqref{eq:Einsbackg} to replace the potential evaluated in the background solution by
\begin{equation}
	V_0=-d(d-1)(A')^2-\frac{1}{2}(\eta_0')^2\ .
\end{equation}
The final form of the equation is
\begin{equation}\label{eq:invariantHeq}
	H''+(d-4)A' H'+\left[e^{-2A}M^2-2(d-2)(A')^2-\frac{1}{d-1}(\eta_0')^2\right]H=0\ .
\end{equation}

To derive the equation for the remaining mode, we introduce
\begin{equation}
	\widehat H_S=\sum_{\alpha=1}^{d-2}e_\mu^\alpha e_\nu^\alpha H^{\mu\nu}, \quad H_L=\frac{k_\mu k_\nu}{k^2}H^{\mu\nu},
\end{equation}
in such a way that the trace is $h=e^{-2A}(H_S+(d-1)\widehat H_S+H_L)$. Then, from the equation for the ghost scalar \eqref{eq:chieq}, the $\mu r$ Einstein equations projected along the momentum, and the $\mu\nu$ Einstein equations projected on $\widehat H_S$ we get the set of equations
\begin{eqnarray}
\label{eq:scalarX} 0&=& X''+dA'X'+e^{-2A}M^2 X+\partial_\eta^2 V_0 X+\frac{1}{2}\eta_0' h'=0,\\
\label{eq:scalarCons} 0&=& H_S'-2A' H_S+(d-1)(\widehat H_S'-2A' \widehat H_S)-e^{2A}\eta_0' X,\\
\label{eq:scalarhHS} 0&=& \widehat H_S''+(d-4) A' \widehat H_S'+\left[4(A')^2+\frac{V_0}{d-1}\right]\widehat H_S+\frac{2}{d-1}e^{2A}\partial_\eta V_0 X+e^{2A}A' h'.
\end{eqnarray}
We subtract \eqref{eq:scalarX}$-\frac{\eta_0'e^{-2A}}{2A'}\eqref{eq:scalarhHS}$ identifying the invariant combination \eqref{eq:zetaeta}. We remove additional terms proportional $\widehat H_S'$ solving for it in \eqref{eq:scalarCons}. The remaining terms depend on $H_S'$, $H_S$ and $\widehat H_S$, but using the equations for the background \eqref{eq:Einsbackg}, \eqref{eq:scalarbackg} one can show that the coefficient of the $\widehat H_S$ term cancels out, so one is left with an equation depending only on invariant quantities
\begin{equation}\label{eq:invariantZeq}
	0=Z_\eta''+dA' Z_\eta'+\left(e^{-2A}M^2+Q_{ZZ}\right)Z_\eta+Q_{ZH}(H_S'-2A' H_S),
\end{equation}
with
\begin{eqnarray}\label{eq:Qcoefs}
	Q_{ZZ}&=& \partial_\eta^2V_0-\frac{1}{d-1}\frac{\eta_0'}{A'}\partial_\eta V_0-\eta_0' e^{2A} Q_{ZH},\\
	Q_{ZH} &=& \frac{e^{-2A}}{d-1}\frac{\eta_0'A''-A'\eta_0''}{(A')^2}.
\end{eqnarray}
Using the background equations of motion \eqref{eq:Einsbackg}, \eqref{eq:scalarbackg} one can also write
\begin{equation}\label{eq:Qcoef2}
	\eta_0'A''-A'\eta_0''=A'\partial_\eta V_0+d(A')^2\eta_0'+\frac{(\eta_0')^3}{2(d-1)}.
\end{equation}

